\chapter{Further Reading}

For the physics of Yang Mills theory and its lattice formulation, Wilson's 1974 paper and Creutz's book \emph{Quarks, Gluons and Lattices} remain the best starting points; the Jaffe and Witten problem description states the mathematical problem precisely. Glimm and Jaffe's \emph{Quantum Physics} is the standard reference for constructive field theory and reflection positivity, and Seiler's lecture notes \emph{Gauge Theories as a Problem of Constructive Quantum Field Theory and Statistical Mechanics} give the Osterwalder and Seiler theory in detail.

For generalized gamma convolutions, Bondesson's monograph is definitive, and Steutel and van Harn's \emph{Infinite Divisibility of Probability Distributions on the Real Line} covers the surrounding theory of infinitely divisible laws. Biane, Pitman and Yor's survey in the Bulletin of the American Mathematical Society is the entry point to the hyperbolic laws and their connection with the zeta function. Karlin's \emph{Total Positivity} covers likelihood ratio orders and P\'olya frequency functions.

For harmonic analysis on compact groups, Br\"ocker and tom Dieck's \emph{Representations of Compact Lie Groups} and Faraut's \emph{Analysis on Lie Groups} are clear and complete. For Brownian motion on groups and two dimensional Yang Mills measures, L\'evy's memoir \emph{Yang Mills Measure on Compact Surfaces} and the papers of Driver and Sengupta are the main sources.

For probabilistic approaches to lattice gauge theory, Chatterjee's survey \emph{Yang Mills for Probabilists} is the natural introduction, and the papers of Cao, Adhikari, Forsstr\"om, Lenells and Viklund develop the finite group case. For the renormalization group, Ba\l aban's papers in Communications in Mathematical Physics from 1984 to 1989 contain the rigorous theory; Kadanoff's notes on Migdal's recursion explain the approximation used in Sections~\ref{ch:mk} and~\ref{ch:verification}.

For Monte Carlo, Sokal's Carg\`ese lectures \emph{Monte Carlo Methods in Statistical Mechanics: Foundations and New Algorithms} are the best account of autocorrelation analysis and cluster algorithms; Gattringer and Lang's \emph{Quantum Chromodynamics on the Lattice} and Montvay and M\"unster's \emph{Quantum Fields on a Lattice} cover heat bath, overrelaxation and glueball measurement. Jerrum's \emph{Counting, Sampling and Integrating} and the book of Levin, Peres and Wilmer, \emph{Markov Chains and Mixing Times}, cover approximate counting and mixing. For Markov random fields, Besag's 1974 paper is still the place to start, and Rue and Held's \emph{Gaussian Markov Random Fields} treats the intrinsic autoregressions. For logarithmic Sobolev inequalities and curvature, Bakry, Gentil and Ledoux, \emph{Analysis and Geometry of Markov Diffusion Operators}, is comprehensive, and Martinelli's Saint Flour lectures cover the equivalence of dynamic and static gaps for spin systems.

For correlation inequalities, Glimm and Jaffe again, and Friedli and Velenik's \emph{Statistical Mechanics of Lattice Systems}. For entire functions with real zeros, Levin's \emph{Distribution of Zeros of Entire Functions} and P\'olya's collected papers. For three dimensional gauge theory, Teper's papers on SU(N) in $2+1$ dimensions and the papers of Karabali, Kim and Nair. For large $N$, Coleman's lecture \emph{$1/N$} in \emph{Aspects of Symmetry}, and for free probability, the lectures of Nica and Speicher.

Our own programme on the Riemann Hypothesis, from which the Thorin tools are drawn, is collected in \emph{Computing the Riemann Hypothesis}, Volumes I and II, and the probabilistic background, with a section on Yang Mills, in \emph{On Hilbert's 6th Problem}; both are on Zenodo, with the digital object identifiers given in the bibliography, and Section~\ref{ch:rhbooks} and Subsection~\ref{ch:h6} describe what this book takes from them.
